\section{Introduction}\label{sec:intro}

The bulk-realization program of this repository asks for an independently
positive system whose state norm is the localized Weil quadratic form
\cite{InverseBulk}. Every candidate tested so far has been asked to reproduce a
\emph{value}: a weight, a pairing, a norm. That has turned out to be a weak
demand. In the one setting where the search was carried out concretely, the set
of reachable states was everything analytic past the unit circle, so an exhibited
weight identity carries no information unless it comes with a reason the
representation could not have produced a different weight \cite{Bottleneck}.

This paper takes up a proposal that asks for something else. A Wilson line, when
deformed, has an expectation value obeying a differential equation fixed by the
variation of its endpoints together with the non-Abelian Stokes theorem. The
shifted family of the program already obeys a differential equation of exactly
that shape, equation (1.7) of the background \cite{Background}:
\[
\frac{\partial V_{\omega,L}}{\partial\omega}=-G_{\omega,L}V_{\omega,L},
\qquad V_{0,L}=I .
\]
The proposal is that these are the same equation, so that the object to search
for is a \emph{deformation} rather than an operator. The attraction is
structural: a generator and an initial condition determine a flow completely, so
matching a flow leaves nothing to tune. It is the right kind of demand to make.

\subsection*{What this paper does}

\textbf{One.} It writes the generator out. The background defined $G_{\omega,L}$
by naming its symbol and inherited the realization silently from the transfer;
Section~\ref{sec:conventions} gives the operator, its domain and its adjoint
status, and Section~\ref{sec:generator} proves that its three factors deliver the
three terms of the target exactly. The resulting statement is compact: the
shifted form is the central form with every off-diagonal kernel multiplied by
$\cosh(\omega(x-y))$ and every prime atom by $\cosh(\omega\log n)$, the diagonal
unchanged (Theorem~\ref{thm:kernel}). The equations of the background are
therefore general --- they carry the whole Weil form, not the archimedean part of
it.

\textbf{Two.} It records a consequence of the causal realization that is easy to
lose. The pole and contact term has \emph{identically zero real part on the
imaginary axis} (Remark~\ref{rem:imaginary}). Reading the generator symbol as a
Fourier multiplier recovers the archimedean term and the prime delays and drops
the contact and poles altogether; they reach the form only through the
right half-plane Laplace realization.

\textbf{Three.} It proves that the shift direction is abelian. Truncation to an
interval is multiplicative on causal kernels (Lemma~\ref{lem:multiplicative}), so
the transfers and their generators lie in one commutative algebra
(Theorem~\ref{thm:abelian}) and the path-ordered solution of the flow equation
collapses to an ordinary exponential. The content of the non-Abelian Stokes
theorem is that path ordering may be traded for surface ordering of the
parallel-transported curvature; for an abelian connection it degenerates, and a
one-parameter family has no surface. \emph{As things stand the mechanism the
proposal names has nothing to act on, and the missing non-commutativity has to be
supplied transverse to $\omega$.} This is the paper's main result and it is a
constraint on the proposal, not a refutation of it.

\textbf{Four.} It identifies what the shift \emph{is}, geometrically. The flow
acts on inputs by the hyperbolic rotation generated by multiplication by $x$
(Proposition~\ref{prop:hyperbolic}), so the shift is a genuine geometric
deformation and its gauge-theoretic analogue is an angle --- a cusp --- rather
than a coupling that can be dialled.

\textbf{Five.} It works out the endpoint direction, which is the one the
compression supplies and the only one left. Three things come out of it
(Section~\ref{sec:endpoint}). The endpoint variation is \emph{rank one}, sitting
at the leading end of the interval with the boundary value of the evolved field
as its coefficient, so the line has a fixed tail and a moving head --- the
endpoint equation of a line, not an analogy for it
(Proposition~\ref{prop:endpoint}). The resulting two-parameter connection is
\emph{flat}, not by hypothesis but because mixed partial derivatives commute
(Proposition~\ref{prop:flat}), so there is no local curvature to extract and the
naive search for one is misconceived. And the connection does not close on the
form domain: endpoint evaluation is not form-bounded, and the smoothing the shift
supplies misses the trace threshold at every admissible shift, by exactly the
borderline (Proposition~\ref{prop:trace}). What survives is the atomic part of
the connection's kernel, whose atoms sit at the prime periods: a flat connection
with punctures, whose monodromy is the one place a non-abelian datum could still
live. That the monodromy is currently one-dimensional, because the residues are
scalars, is what makes the internal index of Remark~\ref{rem:cost} the structural
conclusion rather than a free association.

\subsection*{What this paper does not claim}

No source is constructed. No positivity is proved. Nothing below assumes the
Riemann hypothesis, and no gauge theory is chosen, matched or excluded: the
results are statements about the arithmetic side of a proposed correspondence.
Section~\ref{sec:conditions} proposes three necessary conditions and labels them
as proposals; one of them carries an identification hypothesis that the program
deliberately does not make. Proposition~\ref{prop:trace} is a proof sketch
resting on a smoothing estimate quoted from the background, and
Remark~\ref{rem:marginal} is an observation that explains nothing.

\subsection*{Relation to the companion investigations}

The conventions, the target and the criterion are those of
\cite{InverseBulk,Selection}; Section~\ref{sec:conventions} restates them so that
this paper is self-contained, and it is self-contained apart from the background
\cite{Background}, which it cites by equation number for the shifted family, and
the companions, which it cites for results quoted by number.

Two inherited facts govern the work and are not re-derived here. The first is
that the accumulated exclusions of this program are statements about
\emph{channel decompositions} and about the reachable states of particular
representations; none of them touches a gauge theory, and none of them
constrains a flow, which is not a decomposition \cite{ExclusionMap}. The second
is that a mechanism supplying a margin uniform in $L$ proves something false,
because the relevant contraction is saturated to more than twenty digits
\cite{Selection}. Section~\ref{sec:outlook} argues that the flow formulation is
the one place in the program where that objection may not apply, and states the
question rather than answering it.

Finally, the distinction between lines and loops that motivates the whole
investigation is not a matter of open versus closed geometry. The test functions
carry no boundary conditions; the localization is by support alone. What excluded
the closed benchmark was that its operator had bounded spectrum
\cite{Selection,Bottleneck}, and Section~\ref{sec:conditions} restates that in
the form the deformation picture needs.
